%! From Combinations to Solutions --- Setting Up Ax = b — Unit 2, Lesson 2.0 — Student Packet Cover Sheet
\documentclass[10pt]{article}
\usepackage{linalg-article}
\usepackage{linalg-boxes}
\usepackage{ltablex}
\keepXColumns

\begin{document}

% Full-bleed burgundy banner
\begin{tikzpicture}[remember picture, overlay]
  \fill[burgundy] (current page.north west)
    rectangle ([yshift=-0.9in]current page.north east);
\end{tikzpicture}
\vspace{-0.8in}

\begin{center}
  {\color{white}\textbf{\LARGE Linear Algebra}}\\[4pt]
  {\color{blushmid}\large\textbf{Unit 2 \quad Solving Linear Equations $A\mathbf{x}=\mathbf{b}$}}\\[6pt]
  {\color{white}\textbf{\large Lesson 2.0 \quad From Combinations to Solutions --- Setting Up Ax = b}}
\end{center}

\vspace{0.22in}
\namedateperiod
\vspace{0.18in}

\begin{learningtargetbox}
\small
\begin{enumerate}[leftmargin=1.5em, itemsep=5pt]
  \item \ldots{} read $A\mathbf{x}=\mathbf{b}$ as a \textbf{combination of columns} \emph{and} as a \textbf{system of equations}.
  \item \ldots{} solve a $2\times 2$ system by \textbf{eliminating} a variable, then \textbf{check} by rebuilding $\mathbf{b}$.
  \item \ldots{} explain when a system has \textbf{one}, \textbf{no}, or \textbf{infinitely many} solutions.
\end{enumerate}
\end{learningtargetbox}

\vspace{0.14in}

\begin{tocbox}
\small
\renewcommand{\arraystretch}{1.5}
\begin{tabularx}{\linewidth}{@{} c l X r @{}}
\rowcolor{blush}
\textbf{\#} & \textbf{Component} & \textbf{Description} & \textbf{Score} \\
\midrule
1 & Warm-Up        & Spiral review: solving a system, $A\mathbf{x}$ as columns, reachability & \blank{1.2cm} \\
2 & Guided Notes   & $A\mathbf{x}=\mathbf{b}$ two ways; solving by elimination            & \blank{1.2cm} \\
3 & Group Activity & Blend the Batch --- setting up and solving $A\mathbf{x}=\mathbf{b}$    & \blank{1.2cm} \\
4 & Exit Ticket    & Quick check on setup, elimination, and solvability                   & \blank{1.2cm} \\
5 & Homework       & Three views, more elimination, and a blend application               & \blank{1.2cm} \\
\midrule
  & \multicolumn{2}{r}{\textbf{Total}} & \blank{1.2cm} \\
\end{tabularx}
\end{tocbox}

\vspace{0.10in}

\begin{remindbox}
\small To \textbf{solve} $A\mathbf{x}=\mathbf{b}$ is to find the weights $\mathbf{x}$ that combine
the \textbf{columns} of $A$ into the target $\mathbf{b}$. Each row of $A\mathbf{x}=\mathbf{b}$ is a
\textbf{linear equation}, so solving means finding where those lines meet.
\textbf{Elimination} --- subtracting a multiple of one equation from another to remove a
variable --- is the tool that finds $\mathbf{x}$.
\end{remindbox}

\end{document}
